\section{开源实现与工程实践}
\label{sec:implementation}

从工程实现看，当前生态已经出现了若干可作为"未来模型原生计算架构原型"的项目。
图~\ref{fig:impl_map}将这些系统映射到ICA六层体系结构上。本节按照ICA的层次结构，从服务层（L2）、Agent层（L5）和协议层（L4）三个维度，逐一介绍代表性系统的核心特点及其在ICA中的定位，并提炼出五条可直接指导工程落地的实现建议。

\begin{figure}[htbp]
\centering
\resizebox{0.99\textwidth}{!}{%
\begin{tikzpicture}[
    x=1cm, y=1cm,
    >={Stealth[length=2mm,width=1.3mm]},
    font=\sffamily,
    % --- ICA stack layer band ---
    layerband/.style={
        rounded corners=3pt,
        minimum width=5.0cm,
        minimum height=0.82cm,
        align=center,
        text=white,
        inner xsep=5pt,
        inner ysep=3pt,
        line width=0.5pt,
        draw=white!30,
        font=\small
    },
    % --- Annotation pill ---
    annot/.style={
        rounded corners=2.5pt,
        draw=#1!60,
        fill=#1!8,
        text=#1!85!black,
        font=\scriptsize,
        inner xsep=4pt,
        inner ysep=2.5pt,
        align=center,
        line width=0.5pt
    },
    % --- System box ---
    sysbox/.style={
        rounded corners=2.5pt,
        draw=#1!50,
        fill=#1!6,
        text=black!85,
        font=\scriptsize,
        inner xsep=3.5pt,
        inner ysep=2pt,
        align=center,
        line width=0.5pt
    },
    % --- Recommendation sidebar item ---
    recbox/.style={
        rounded corners=3pt,
        draw=ink!30,
        fill=ink!4,
        text=ink,
        font=\scriptsize,
        inner xsep=4pt,
        inner ysep=3pt,
        align=left,
        line width=0.5pt,
        text width=3.8cm
    },
    % --- Connector arrow ---
    conn/.style={
        -{Stealth[length=1.6mm,width=1.1mm]},
        draw=#1!55,
        line width=0.6pt
    },
    % --- Section header ---
    colhead/.style={
        font=\footnotesize\bfseries,
        text=ink
    },
    layernum/.style={
        font=\scriptsize\bfseries,
        text=softgray
    },
]

% =========================================================================
% Color definitions (ensure they exist)
% =========================================================================
\providecolor{classical}{HTML}{3D6FB6}
\providecolor{modelnative}{HTML}{D9792B}
\providecolor{ink}{HTML}{22303C}
\providecolor{softgray}{HTML}{6E7B87}
\providecolor{bgline}{HTML}{D8DEE6}

% Layer band colors (consistent with fig_dual_stack)
\providecolor{c1}{HTML}{234C84}
\providecolor{c2}{HTML}{2F5F9E}
\providecolor{c3}{HTML}{3E73B2}
\providecolor{m4}{HTML}{E38431}
\providecolor{m5}{HTML}{EE9D50}
\providecolor{g6}{HTML}{7B8A9E}

% =========================================================================
% Coordinates
% =========================================================================
% Stack center x=0, layers from bottom (L1) to top (L6)
% y positions: L1=0, L2=1.2, L3=2.4, L4=4.2, L5=5.4, L6=6.6
% (gap at L3/L4 for dashed line)

% =========================================================================
% Background: Probabilistic Execution Plane (L1--L3)
% =========================================================================
\fill[classical!8, rounded corners=8pt]
    (-2.8,-0.55) rectangle (2.8,3.15);

% =========================================================================
% Background: Deterministic Control Plane (L4--L5)
% =========================================================================
\fill[modelnative!8, rounded corners=8pt]
    (-2.8,3.55) rectangle (2.8,6.1);

% =========================================================================
% Plane labels (small, inside background areas)
% =========================================================================
\node[font=\tiny\bfseries, text=classical!60!black, anchor=north west]
    at (-2.65,3.05) {Probabilistic Execution Plane};
\node[font=\tiny\bfseries, text=modelnative!60!black, anchor=north west]
    at (-2.65,6.0) {Deterministic Control Plane};

% =========================================================================
% ICA Layer Bands (central stack)
% =========================================================================
% L1 -- Physical Execution
\node[layerband, fill=c1] (L1) at (0,0)
    {\textbf{L1}~~Physical Execution};

% L2 -- Inference Serving
\node[layerband, fill=c2] (L2) at (0,1.15)
    {\textbf{L2}~~Inference Serving};

% L3 -- Context Management
\node[layerband, fill=c3] (L3) at (0,2.3)
    {\textbf{L3}~~Context Management};

% ---- Dashed boundary at L3/L4 ----
\draw[dashed, black!40, line width=0.9pt]
    (-3.5,3.35) -- (3.5,3.35);
\node[font=\tiny\itshape, text=black!40, anchor=west]
    at (-3.4,3.35) {L3/L4 interface};

% L4 -- Semantic Interface
\node[layerband, fill=m4] (L4) at (0,4.1)
    {\textbf{L4}~~Semantic Interface};

% L5 -- Orchestration
\node[layerband, fill=m5] (L5) at (0,5.2)
    {\textbf{L5}~~Orchestration};

% L6 -- Application
\node[layerband, fill=g6] (L6) at (0,6.3)
    {\textbf{L6}~~Application};

% =========================================================================
% Layer number labels (far left)
% =========================================================================
\foreach \y/\lab in {0/L1, 1.15/L2, 2.3/L3, 4.1/L4, 5.2/L5, 6.3/L6}{
    \node[layernum] at (-4.3,\y) {\lab};
    \draw[bgline, line width=0.4pt] (-4.0,\y) -- (-2.75,\y);
}

% =========================================================================
% Vertical stack arrows
% =========================================================================
\foreach \a/\b in {L1/L2, L2/L3, L3/L4, L4/L5, L5/L6}{
    \draw[-{Stealth[length=1.5mm]}, black!25, line width=0.5pt] (\a) -- (\b);
}

% =========================================================================
% COLUMN 1 -- Serving Layer (left, x ~ -6.5 to -3.5)
% =========================================================================
\node[colhead] at (-5.4,7.3) {Serving Layer};

% L1 systems
\node[sysbox=classical, text width=2.6cm] (TGI) at (-6.5,0.3)
    {\textbf{TGI}\\[-1pt]{\tiny Unified multi-model interface}};
\node[sysbox=classical, text width=2.6cm] (TRT) at (-4.3,0.3)
    {\textbf{TensorRT-LLM}\\[-1pt]{\tiny FP8, Tensor Core}};

% Annotation: kernel fusion
\node[annot=classical, text width=2.0cm] (annkf) at (-5.4,-0.65)
    {kernel fusion};
\draw[conn=classical] (annkf.north) -- (-5.4,-0.05);

% Connect to L1
\draw[conn=classical] (TGI.east) -- ([xshift=-4pt]L1.west);
\draw[conn=classical] (TRT.east) -- ([xshift=-4pt]L1.west);

% L2 systems
\node[sysbox=classical, text width=2.6cm] (vLLM) at (-6.5,1.45)
    {\textbf{vLLM}\\[-1pt]{\tiny PagedAttention}};
\node[sysbox=classical, text width=2.6cm] (SGLang) at (-4.3,1.45)
    {\textbf{SGLang}\\[-1pt]{\tiny RadixAttention}};

% Connect to L2
\draw[conn=classical] (vLLM.east) -- ([xshift=-4pt]L2.west);
\draw[conn=classical] (SGLang.east) -- ([xshift=-4pt]L2.west);

% =========================================================================
% COLUMN 2 -- Protocol Layer (center-top, x ~ -1.5 to 1.5)
% =========================================================================
\node[colhead] at (0,7.3) {Protocol Layer};

% MCP -- vertical bus (spans L4-L5)
\node[sysbox=modelnative, text width=2.2cm, minimum height=1.6cm] (MCP) at (-1.6,4.65)
    {\textbf{MCP}\\[2pt]{\tiny JSON-RPC\\tool server bus}};

% MCP annotation
\node[annot=modelnative, text width=1.6cm] (annmcp) at (-1.6,3.55)
    {tool discovery};
\draw[conn=modelnative] (annmcp.north) -- (MCP.south);

% Connect MCP to L4 and L5
\draw[conn=modelnative] (MCP.west) -- ([xshift=2pt]L4.east);
\draw[conn=modelnative] ([yshift=4pt]MCP.west) -- ([xshift=2pt]L5.east);

% A2A -- horizontal interconnect (spans L5-L6)
\node[sysbox=modelnative, text width=2.2cm, minimum height=1.1cm] (A2A) at (1.6,5.7)
    {\textbf{A2A}\\[2pt]{\tiny agent-to-agent}};

% A2A annotation
\node[annot=modelnative, text width=1.8cm] (anna2a) at (1.6,6.8)
    {agent communication};
\draw[conn=modelnative] (anna2a.south) -- (A2A.north);

% Connect A2A to L5 and L6
\draw[conn=modelnative] (A2A.west) -- ([xshift=6pt]L5.east);
\draw[conn=modelnative] ([yshift=-3pt]A2A.west) -- ([xshift=4pt]L6.east);

% =========================================================================
% COLUMN 3 -- Agent Layer (right, x ~ 4.5 to 8.0)
% =========================================================================
\node[colhead] at (6.5,7.3) {Agent Layer};

% L5: Industrial agents
\node[sysbox=modelnative, text width=1.7cm] (Codex) at (4.7,5.5)
    {\textbf{Codex}};
\node[sysbox=modelnative, text width=1.7cm] (Claude) at (6.5,5.5)
    {\textbf{Claude Code}};
\node[sysbox=modelnative, text width=1.7cm] (OH) at (8.3,5.5)
    {\textbf{OpenHands}};

% Connect to L5
\draw[conn=modelnative] (Codex.west) -- ([xshift=8pt]L5.east);
\draw[conn=modelnative] (Claude.west) -- ([xshift=12pt]L5.east);
\draw[conn=modelnative] (OH.west) -- ([xshift=16pt]L5.east);

% L4-L5 spanning: Agent OS
\node[sysbox=modelnative, text width=3.8cm, minimum height=0.8cm] (ArbOS) at (6.5,4.35)
    {\textbf{ArbiterOS / AgentOS}\\[-1pt]{\tiny governance $\cdot$ reasoning kernel $\cdot$ sandboxing}};

% Connect ArbiterOS to L4 and L5
\draw[conn=modelnative] (ArbOS.west) -- ([xshift=10pt]L4.east);
\draw[conn=modelnative] ([yshift=3pt]ArbOS.west) -- ([xshift=14pt]L5.east);

% =========================================================================
% RECOMMENDATIONS SIDEBAR (far right)
% =========================================================================
\node[font=\footnotesize\bfseries, text=ink, anchor=south] at (13.2,7.3)
    {Implementation};
\node[font=\footnotesize\bfseries, text=ink, anchor=north] at (13.2,7.2)
    {Recommendations};

% Rec 1: Separate control plane from data plane
\node[recbox] (R1) at (13.2,6.15)
    {\textbf{R1.} Separate control \& data planes\\[-1pt]{\tiny Codex approvals, Claude Code hooks}};

% Rec 2: Separate short-term context from long-term memory
\node[recbox] (R2) at (13.2,4.85)
    {\textbf{R2.} Separate hot context \& cold memory\\[-1pt]{\tiny MemGPT virtual context, L3 tiering}};

% Rec 3: Version every behavior-affecting object
\node[recbox] (R3) at (13.2,3.55)
    {\textbf{R3.} Version every behavior object\\[-1pt]{\tiny MCP schemas, tool ABI versioning}};

% Rec 4: Enforce resource quotas
\node[recbox] (R4) at (13.2,2.25)
    {\textbf{R4.} Enforce resource quotas\\[-1pt]{\tiny cgroup-style budgets, Axiom 5 \& 6}};

% Rec 5: Design failure recovery
\node[recbox] (R5) at (13.2,0.95)
    {\textbf{R5.} Design failure recovery\\[-1pt]{\tiny Checkpoint/retry, distributed DB patterns}};

% =========================================================================
% Connector lines: Recommendations to specific systems/layers
% =========================================================================
% R1 -> L5 (Codex, Claude Code, ArbiterOS)
\draw[conn=ink, dashed] (R1.west) -- ++(-1.2,0) |- ([xshift=2pt]Claude.east);

% R2 -> L3 (context management)
\draw[conn=ink, dashed] (R2.west) -- ++(-0.8,0) |- ([xshift=14pt]L3.east);

% R3 -> MCP / L4
\draw[conn=ink, dashed] (R3.west) -- ++(-1.0,0) |- ([xshift=-2pt]MCP.east);

% R4 -> spans L4-L5
\draw[conn=ink, dashed] (R4.west) -- ++(-0.8,0) |- ([xshift=18pt]L4.east);

% R5 -> spans L2-L5
\draw[conn=ink, dashed] (R5.west) -- ++(-0.6,0) |- ([xshift=12pt]L2.east);

% =========================================================================
% Decorative: vertical brace for recommendations
% =========================================================================
\draw[decorate, decoration={brace, amplitude=5pt, mirror},
      line width=0.7pt, draw=ink!40]
    (15.3,6.7) -- (15.3,0.45)
    node[midway, right=8pt, font=\tiny\bfseries, text=ink!60, align=center]
    {Five\\recs};

% =========================================================================
% Subtle footnote
% =========================================================================
\node[font=\tiny, text=softgray, anchor=west] at (-6.5,-1.1)
    {Open-source ecosystem mapped to the six-layer ICA architecture; dashed lines connect recommendations to target layers.};

\end{tikzpicture}%
}
\caption{Open-source ecosystem mapped to the ICA architecture.  The central stack shows
  the six ICA layers with dual-plane shading (blue = probabilistic execution plane L1--L3;
  orange = deterministic control plane L4--L5).  Three annotation columns identify
  representative systems at the serving, protocol, and agent layers.  The right sidebar
  connects the five implementation recommendations to the specific layers and systems
  they target.}
\label{fig:impl_map}
\end{figure}


\subsection{服务层（L2）：推理服务与KV缓存管理}
\label{sec:impl-serving}

推理服务层对应ICA的L2层，负责模型权重加载、KV缓存管理、连续批处理和预填充/解码调度。
目前这一层已形成较为成熟的工程生态。

\textbf{vLLM}以PagedAttention为核心，将KV缓存组织为固定大小的页（block），通过block table映射到非连续物理显存空间，实现了类似操作系统虚拟内存的显存管理\cite{kwon2023pagedattention,vllm_docs_2026}。
其自动前缀缓存（Automatic Prefix Caching）功能进一步支持了跨请求的共享前缀复用，相当于L2缓存中的共享只读代码页\cite{vllm_prefix_2026}。

\textbf{SGLang}强调结构化程序执行与RadixAttention\cite{sglang2024,sglang_docs_2026}。
RadixAttention用基数树（radix tree）索引所有历史请求的token前缀，自动发现和复用最大公共前缀的KV缓存，在多轮对话和批量推理场景下展现出更高的缓存命中率。
实测表明，SGLang在某些工作负载下吞吐量比vLLM高出约29\%\cite{sglang2024}。

\textbf{TGI}（Text Generation Inference）与\textbf{TensorRT-LLM}分别由Hugging Face和NVIDIA提供，展示了面向生产环境的部署能力\cite{tgi_docs_2026,tensorrtllm_docs_2026}。
前者侧重多模型统一接口和易用性，后者利用NVIDIA GPU的深度优化（如FP8量化、Tensor Core加速）追求极致吞吐。
这三类系统（vLLM、SGLang、TGI/TensorRT-LLM）分别代表了"类虚拟内存管理""结构化缓存复用"和"硬件深度优化"三种工程路线，共同构成了L2层的实现谱系。

\subsection{Agent层（L5）：Agent运行时与操作系统}
\label{sec:impl-agent}

Agent层对应ICA的L5层，负责任务规划、Agent调度、权限审批、失败恢复和状态提交。
这一层的开源生态最为活跃，方向也最为多样，可分为三个子方向。

\subsubsection{通用Agent}

\textbf{OpenHands}提供了一个开放的Agent平台，支持多种LLM后端、沙箱执行和可扩展的工具接口\cite{openhands_paper_2025}；
\textbf{SWE-agent}专注于软件工程任务，通过Agent-Computer Interface（ACI）将代码编辑、测试和调试统一为结构化操作\cite{sweagent2024}；
\textbf{AutoGPT}从自主持续工作流出发，探索了Agent的自动目标分解和长期执行能力\cite{autogpt_platform_2026}；
\textbf{Letta}（原MemGPT团队）则聚焦于"有状态Agent"，将长期记忆作为Agent的一等公民来管理\cite{letta_stateful_2026}。

\subsubsection{Agent操作系统}

多个项目已开始将Agent运行时本身视为操作系统。\textbf{ArbiterOS}引入了治理层概念，提出"governor治理概率CPU"的分层架构\cite{arbiteros2025}；
\textbf{AgentOS}提出了受结构化操作系统逻辑约束的"reasoning kernel"\cite{agentos2025}；
\textbf{UFO$^2$}展示了桌面级AgentOS的可行性，将HostAgent、AppAgent和GUI--API操作层组织为层次结构\cite{ufo2_2025}；
\textbf{Aura}则在移动端探索了安全优先的Agent架构，包含System Agent、沙箱化的App Agents和Agent Kernel\cite{aura2025}。

\subsubsection{工业Agent系统}

两个工业系统特别值得关注，因为它们明确采用了操作系统构建的惯例。\textbf{OpenAI Codex}明确将子Agent（subagents）、技能（skills）、MCP服务器接入、沙箱执行和审批策略（approval policies）作为一等公民\cite{openai_codex_overview_2026,openai_codex_skills_2026,openai_codex_sandbox_2026,openai_codex_approvals_2026}；
\textbf{Anthropic Claude Code}提供子Agent、技能、hooks、MCP连接和项目记忆文件，并默认采用只读权限与审批机制\cite{anthropic_claudecode_overview_2026,anthropic_claudecode_security_2026}。
这两个系统说明，工业界已经开始把Agent运行时当作操作系统来构建。

\subsection{协议层（L4）：标准化工具连接与Agent互联}
\label{sec:impl-protocol}

协议层对应ICA的L4层，负责统一外部能力的访问接口。
这一层正在从"点对点工具调用"向"标准化总线协议"转变。

\textbf{MCP}（Model Context Protocol）提供了标准化的工具连接协议\cite{mcp_intro_2026}。
截至2025年底，Anthropic报告已有超过10{,}000个活跃的公开MCP服务器\cite{mcpadoption2026}。
MCP将连接AI应用与外部数据源、工具、工作流统一为JSON-RPC规范，扮演类似PCIe总线或USB在传统计算机中的角色。

\textbf{A2A}（Agent-to-Agent Protocol）由Google于2025年4月推出，定位为Agent间通信标准\cite{google_a2a_2025,agentprotocols2025}。
它支持能力发现（capability discovery）、任务生命周期管理和多模态协同，已被捐赠给Linux基金会进行中立治理\cite{a2alinux2025}。
A2A扮演类似网络协议栈中TCP/IP的角色，与MCP互补：MCP是Agent到工具的\textbf{纵向总线}，A2A是Agent到Agent的\textbf{横向互联}。

\subsection{五条实现建议}
\label{sec:impl-recommendations}

基于上述工程实践和本文提出的六条设计公理，本文给出五条可直接指导工程落地的实现建议。图~\ref{fig:impl_recommendations}将每条建议映射到其主要ICA层和对应公理。

% ---------------------------------------------------------------------------
% Five Implementation Recommendations × ICA Layer mapping
% fig:impl_recommendations
% ---------------------------------------------------------------------------
\begin{figure}[htbp]
\centering
\providecolor{r1Color}{HTML}{4A78A8}
\providecolor{r2Color}{HTML}{5C9ED4}
\providecolor{r3Color}{HTML}{E07B39}
\providecolor{r4Color}{HTML}{C44E52}
\providecolor{r5Color}{HTML}{3A8A3A}
\providecolor{dotFull}{HTML}{333333}
\providecolor{dotPart}{HTML}{999999}
%
\resizebox{0.8\textwidth}{!}{%
\begin{tikzpicture}[
    font=\small, >=Latex,
    recbox/.style={draw=#1!65!black, fill=#1!14, rounded corners=5pt,
        minimum width=3.4cm, minimum height=1.2cm,
        align=center, font=\small\bfseries, text=#1!80!black, line width=0.7pt,
        inner sep=5pt},
    layercell/.style={draw=#1!50, fill=#1!12, rounded corners=3pt,
        minimum width=1.55cm, minimum height=0.75cm,
        align=center, font=\scriptsize, text=#1!75!black, line width=0.4pt},
    nabox/.style={fill=gray!8, draw=gray!25, rounded corners=3pt,
        minimum width=1.55cm, minimum height=0.75cm,
        align=center, font=\scriptsize, text=gray!50, line width=0.3pt},
    axiomtag/.style={draw=#1!55, fill=#1!10, rounded corners=2pt,
        font=\tiny\bfseries, text=#1!75!black, inner xsep=3pt, inner ysep=2pt,
        line width=0.4pt},
    hdr/.style={font=\small\bfseries, text=black!65},
]

% ---- Column header positions ----
\def\xR{-0.2}   % Recommendation col
\def\xLone{2.8} \def\xLtwo{4.5} \def\xLthree{6.2}
\def\xLfour{7.9} \def\xLfive{9.6} \def\xLsix{11.3}
\def\xAxiom{12.8}

% ---- Headers ----
\node[hdr] at (\xR, 6.5)     {Recommendation};
\node[hdr, text=blue!60!black]   at (\xLone,   6.5) {L1};
\node[hdr, text=blue!60!black]   at (\xLtwo,   6.5) {L2};
\node[hdr, text=blue!60!black]   at (\xLthree, 6.5) {L3};
\node[hdr, text=orange!60!black] at (\xLfour,  6.5) {L4};
\node[hdr, text=orange!60!black] at (\xLfive,  6.5) {L5};
\node[hdr, text=orange!60!black] at (\xLsix,   6.5) {L6};
\node[hdr, text=red!60!black]    at (\xAxiom,  6.5) {Axiom};

% ---- R1: Separate control & data planes ----
\def\yRa{5.1}
\node[recbox=r1Color, text width=3.0cm] (r1) at (\xR, \yRa)
    {R1\\Separate control\\$\&$ data planes};
\node[nabox]          at (\xLone,   \yRa) {---};
\node[nabox]          at (\xLtwo,   \yRa) {---};
\node[nabox]          at (\xLthree, \yRa) {---};
\node[layercell=r1Color] at (\xLfour,  \yRa) {Key};
\node[layercell=r1Color] at (\xLfive,  \yRa) {Key};
\node[nabox]          at (\xLsix,   \yRa) {---};
\node[axiomtag=r1Color]  at (\xAxiom,  \yRa) {A3, A6};

% ---- R2: Separate hot context & cold memory ----
\def\yRb{3.4}
\node[recbox=r2Color, text width=3.0cm] (r2) at (\xR, \yRb)
    {R2\\Hot context /\\cold memory};
\node[nabox]          at (\xLone,   \yRb) {---};
\node[layercell=r2Color] at (\xLtwo,   \yRb) {Partial};
\node[layercell=r2Color] at (\xLthree, \yRb) {Key};
\node[nabox]          at (\xLfour,  \yRb) {---};
\node[nabox]          at (\xLfive,  \yRb) {---};
\node[nabox]          at (\xLsix,   \yRb) {---};
\node[axiomtag=r2Color]  at (\xAxiom,  \yRb) {A1, A4};

% ---- R3: Version every behavior object ----
\def\yRc{1.7}
\node[recbox=r3Color, text width=3.0cm] (r3) at (\xR, \yRc)
    {R3\\Version every\\behavior object};
\node[nabox]          at (\xLone,   \yRc) {---};
\node[nabox]          at (\xLtwo,   \yRc) {---};
\node[layercell=r3Color] at (\xLthree, \yRc) {Partial};
\node[layercell=r3Color] at (\xLfour,  \yRc) {Key};
\node[layercell=r3Color] at (\xLfive,  \yRc) {Partial};
\node[nabox]          at (\xLsix,   \yRc) {---};
\node[axiomtag=r3Color]  at (\xAxiom,  \yRc) {A2};

% ---- R4: Enforce resource quotas ----
\def\yRd{0.0}
\node[recbox=r4Color, text width=3.0cm] (r4) at (\xR, \yRd)
    {R4\\Enforce resource\\quotas};
\node[nabox]          at (\xLone,   \yRd) {---};
\node[layercell=r4Color] at (\xLtwo,   \yRd) {Partial};
\node[layercell=r4Color] at (\xLthree, \yRd) {Partial};
\node[nabox]          at (\xLfour,  \yRd) {---};
\node[layercell=r4Color] at (\xLfive,  \yRd) {Key};
\node[nabox]          at (\xLsix,   \yRd) {---};
\node[axiomtag=r4Color]  at (\xAxiom,  \yRd) {A5, A6};

% ---- R5: Design failure recovery ----
\def\yRe{-1.7}
\node[recbox=r5Color, text width=3.0cm] (r5) at (\xR, \yRe)
    {R5\\Design failure\\recovery};
\node[nabox]          at (\xLone,   \yRe) {---};
\node[layercell=r5Color] at (\xLtwo,   \yRe) {Partial};
\node[layercell=r5Color] at (\xLthree, \yRe) {Partial};
\node[nabox]          at (\xLfour,  \yRe) {---};
\node[layercell=r5Color] at (\xLfive,  \yRe) {Key};
\node[layercell=r5Color] at (\xLsix,  \yRe) {Partial};
\node[axiomtag=r5Color]  at (\xAxiom,  \yRe) {A3, A6};

% ---- Legend (single row, centred under table) ----
\def\legy{-3.1}
% Diagram spans x: -0.2 to ~14.0, mid ~6.9
% Three items each ~3.8cm wide, total ~11.4cm, start at ~1.2
\node[layercell=dotFull, minimum width=1.3cm, text=dotFull!70] at (0.6, \legy) {Key};
\node[font=\scriptsize, anchor=west] at (1.45, \legy) {Primary layer};
\node[layercell=dotPart, minimum width=1.3cm, text=dotPart!70] at (5.1, \legy) {Partial};
\node[font=\scriptsize, anchor=west] at (5.95, \legy) {Supporting role};
\node[nabox, minimum width=1.3cm] at (9.6, \legy) {---};
\node[font=\scriptsize, anchor=west] at (10.45, \legy) {Not applicable};

\end{tikzpicture}}%
\caption{Mapping of the five implementation recommendations (R1--R5) onto ICA layers (L1--L6) and their governing design axioms. \textbf{Key} cells indicate the primary ICA layer at which the recommendation must be enforced; \textbf{Partial} cells indicate supporting roles. R1 (control/data plane separation) targets L4--L5; R2 (context tiering) targets L3; R3 (versioning) targets L4; R4 (resource quotas) and R5 (failure recovery) both center on L5 with support from adjacent layers.}
\label{fig:impl_recommendations}
\end{figure}


\subsubsection{控制面与数据面分离}

将确定性控制逻辑（权限审批、日志审计、失败恢复）从概率式推理路径中剥离，分别对应双平面架构中的确定性控制平面和概率执行平面。
例如，Codex的approval policies和Claude Code的hooks就是这一分离的具体实现\cite{openai_codex_approvals_2026,anthropic_claudecode_security_2026}。

\subsubsection{短期上下文与长期记忆分离}

将高频使用的"热上下文"（当前对话窗口）与低频访问的"冷记忆"（历史摘要、知识库索引）分到不同的存储介质和管理策略中。
这对应ICA L3层的冷热分层设计，MemGPT的虚拟上下文管理就是这一原则的典型实践\cite{memgpt2023}。
Augment Code的企业案例——CTO预估4--8个月的项目在2周内完成\cite{anthropic2026agenticreport}——展示了有效的上下文管理对加速开发者入职和项目交付的实际影响。

\subsubsection{所有可影响行为的对象都要版本化}

提示模板、工具schema、技能包、记忆文件、MCP服务器能力描述都应附带版本号和变更日志\cite{mcp_intro_2026,agentprotocols2025}。
传统计算机通过ISA和ABI的稳定性保证向后兼容；模型原生计算系统同样需要这种"接口契约"来应对快速演化。

\subsubsection{像OS一样实施资源配额}

对每个Agent或任务设定显式的资源预算：上下文窗口上限、工具调用次数上限、推理时间上限和成本上限\cite{linux_cgroup_2026}。
这对应公理5（最小权限公理）和公理6（可观测性公理）在资源管理层面的具体化。

\subsubsection{像分布式系统一样设计故障恢复}

Agent执行的长任务可能因模型超时、工具失败或上下文溢出而中断。
系统应像分布式数据库一样提供检查点（checkpoint）、重试（retry）和回滚（rollback）机制，确保副作用操作可审计、可撤销\cite{raft2014,spanner2012}。
Rakuten的Claude Code自主完成了vLLM项目中1250万行代码上的激活向量提取任务（7小时，99.9\%数值精度）\cite{anthropic2026agenticreport}，表明具备隐式检查点和重试机制的长时运行Agent已经能够处理工业级任务。

%% ============================================================
